Coral reefs are an existential habitat for a wide range of marine species, offering shelter, 
food and breeding grounds. However, they are highly sensitive to environmental changes, and 
are increasingly threatened by climate change and human activities. When preservation is not 
possible, restoration is seen as a promising alternative to regenerate the biome and preserve 
the marine biodiversity. Current approaches and techniques for coral restoration can't avoid
an intrinsic risk for human divers. For these reasons, the use of autonomous underwater vehicles
(AUVs) has been proposed as a solution to overcome the limitations of human divers and to increase
safety and efficiency in the search task. Moreover, given the unknown and complex nature of 
the marine environment, MARL algorithms have been increasingly used in these kind of tasks.
This work aims to test and evaluate the use of a suitable MARL algorithm in a phisically realistic 
simulated environment, where the agents are required to navigate and search for suitable target zones
defined by a set of environmental characteristics. 
The results show that the proposed MARL algorithm is able to learn an optimal policy to navigate
an unknown environment and find the target zones, outperforming a random policy. Moreover, 
the use of independent agents lead to better performance overall, against the use of a centralized 
critic or communication at execution time, even though in specific scenarios their use can be beneficial 
and lead to higher efficiency.